% !PRINTED 214
\weberpage{214}
%<<CONT
group $\mathfrak{G}$ of order $\mu$ is certainly contained in $\mathfrak{A}$. From this there results

\weq{16}{\mu = \nu\varphi(n)}

and at the same time the identity of $\mathfrak{G}$ with $\mathfrak{A}$.

But it also follows from this that $F(t)$ is equal to the product of all the factors $F(t,\varrho^m)$, so that

\weq{17}{F(t) = \overset{m}{\varPi} F(t, \varrho^m),}

and, since $F(t) = 0$ has no repeated roots, that $F(r,\varrho)$ does indeed vanish, but that none of the other factors $F(r,\varrho^m)$ does. The equations

\weq{18}{F(r,t) = 0, \quad \varPhi(t) = 0}

have therefore only the single root $t = \varrho$ in common with one another, and by seeking out their greatest common divisor \gesperrt{one finds $\varrho$ expressed rationally by $r$, that is, by the roots of the division equation.}

These expressions for $\varrho$ do not change their value when a substitution of the group $\mathfrak{B}$ is applied to $r$, while [according to (12)] $\varrho$ passes over into a power of $\varrho$ under a substitution of the group $\mathfrak{A}$.

The \gesperrt{Abelian} relations show that under the substitution contained in $\mathfrak{A}$

\weq{19}{\begin{pmatrix} m, & 0 \\ 0, & 1 \end{pmatrix}}

$\varrho$ passes over into $\varrho^m$; for if we set

\[
e^{\frac{8\pi i}{n}} = \varrho^\lambda,
\]

then one of the \gesperrt{Abelian} relations reads

\weq{20}{\overset{\nu}{\Sigma} \varrho^{\lambda\nu\nu'} x_{\nu,\nu'} = 0,}

to which, when the expression in terms of $r$ is put in place of $\varrho$, \gesperrt{all} substitutions of $\mathfrak{A}$, and hence (19) as well, are applicable. But by § 62, (12), equation (20) remains correct under this substitution \gesperrt{only} if $\varrho$ passes over into $\varrho^m$.

Since $\varrho$ is not changed by the substitutions in $\mathfrak{B}$, and since the group $\mathfrak{A}$ arises from $\mathfrak{B}$ by composition with (19), it follows that under any substitution in $\mathfrak{A}$

\[
\begin{pmatrix} a, & b \\ c, & \partial \end{pmatrix}
\]
%CONT>>
